\section{Preliminaries: GRPO and Rollout Staleness}
\label{sec:preliminaries}
We briefly review Group Relative Policy Optimization
(GRPO)~\citep{shao2024deepseekmathpushinglimitsmathematical} and define rollout
staleness. Let $\beta$ denote the behavior policy that generates rollouts and
$\pi_\theta$ the policy being optimized.

\paragraph{GRPO objective.} For each prompt $x\sim\mathcal D$, the behavior policy $\beta$ samples a group
of $G$ responses $\{y_i\}_{i=1}^G$, where $y_i=(a_{i,1},\ldots,a_{i,T_i})$.
Each response receives a scalar reward $R_i$. GRPO estimates advantages by
normalizing rewards within the group:
\begin{equation}
    A_i
    =
    \frac{
        R_i - \operatorname{mean}(\{R_1,\ldots,R_G\})
    }{
        \operatorname{std}(\{R_1,\ldots,R_G\})
    } .
    \label{eq:grpo-adv}
\end{equation}
The same response-level advantage $A_i$ is assigned to all tokens in response
$y_i$.

For token $a_{i,t}$, with prefix state $s_{i,t}=(x,a_{i,<t})$, define the
token-level importance ratio
$
    \rho_{i,t}(\theta)
    =
    \frac{
        \pi_\theta(a_{i,t}\mid s_{i,t})
    }{
        \beta(a_{i,t}\mid s_{i,t})
    } .
    \label{eq:token-ratio-prelim}
$
GRPO optimizes
\begin{equation}
\begin{aligned}
    \mathcal J_{\mathrm{GRPO}}(\theta)
    =
    \mathbb E_{x\sim\mathcal D,\; \{y_i\}\sim\beta}
    \left[
    \frac{1}{G}\sum_{i=1}^G
    \frac{1}{T_i}\sum_{t=1}^{T_i}
    \min\left(
        \rho_{i,t}(\theta) A_i,\,
        \operatorname{clip}
        \big(\rho_{i,t}(\theta),1-\epsilon,1+\epsilon\big) A_i
    \right)
    \right].
\end{aligned}
\label{eq:grpo-obj}
\end{equation}
The implemented training loss is $-\mathcal J_{\mathrm{GRPO}}$. 
A KL term $-\lambda_{\mathrm{KL}} D_{\mathrm{KL}}\!\left(
\pi_\theta \,\|\, \pi_{\mathrm{ref}}
\right)$ is optionally applied.

\paragraph{Rollout staleness $\mu$.} GRPO training alternates between rollout generation and policy optimization.
At the beginning of each cycle, the current policy is frozen as the behavior
policy $\beta$, which generates a fixed rollout batch used for several policy
updates before the next refresh.

Let $B_{\mathrm{train}}$ denote the number of prompt groups generated in one
rollout phase, where each prompt group contains $G$ sampled responses. Let
$B_{\mathrm{mini}}$ denote the number of prompt groups consumed by one policy
update. Therefore, a
single rollout generation phase supports $\mu = {B_{\rm train}}/{B_{\rm mini}}$ model updates, which we refer to as the \textit{rollout staleness}.

During these $\mu$ updates, $\theta$ changes while the behavior policy $\beta$
and the stored rollout log-probabilities remain fixed. Thus later updates within
the same cycle optimize against increasingly stale data. A strictly on-policy
algorithm corresponds to $\mu=1$, where every update uses freshly generated
rollouts. 

GRPO is therefore a \emph{near on-policy} algorithm, with $\mu$ typically set to a small value (e.g., $\mu\,$=$\,4$~\citep{zheng2025act,wang2025reinforcement,yu2025dapoopensourcellmreinforcement,chen2025minimax,zheng2025group}) to control policy drift under PPO-style clipped updates. 
While it is generally understood that a large $\mu$ degrades performance,\textbf{ we show in this work that substantially larger values of \boldmath $\mu$ are feasible}, achieving comparable or improved performance while accelerating training.

Throughout the paper, one \emph{model update} refers to one optimizer update on
one minibatch, and \emph{batch size} refers to $B_{\mathrm{mini}}$, the number of prompt groups for
one optimizer update, unless stated
otherwise.